\section{Proposed Resident-Streaming Roofline}
\label{sec:model}
\subsection{Workload Demands and Hardware Throughputs}
Following the operand roles in Fig.~\ref{fig:operand-roles}(b), the model analyzes a \emph{workload window}, which specifies the input work and the resident state at its start.

Within the window, the streaming demand $Q_S$ is the number of streaming input bytes served at the chosen logical input boundary, and the resident demand $Q_R$ is the number of logical bytes actually written to load, update, or reload resident state. An initial load counts only if it occurs inside the window. The resident operand may be a weight or, in attention, the KV cache of earlier tokens.

Each demand is served by its own throughput. The streaming throughput $\rho$ is the maximum rate at which the hardware serves logical input bytes while completing the specified calculation and outputs. The resident throughput $\tau$ is the maximum rate at which logical resident bytes become ready for calculation under a given write mode. Each demand and its throughput share the same configuration and precision. We define
\begin{equation}
  \SI=\frac{Q_S}{Q_R},\qquad P=\frac{Q_S}{T}\quad[\Byte/\mathrm{s}],
  \label{eq:definitions}
\end{equation}
where $T$ is the window completion time, $P$ is the average streaming-input throughput over the window, and $\SI$ is the streaming demand per byte of resident write, with $\SI=\infty$ for $Q_S>0$ and $Q_R=0$.

Demands count logical bytes; the physical steps that serve them, such as bit-serial cycles, repeated conversions, or program and verify steps, enter the service time of the corresponding path. Maintenance that preserves existing state, such as refresh, also consumes service time but is not a new logical update. A repeated input is counted again only when it re-crosses the chosen boundary, and each demand is paired with the throughput defined at that boundary.

\paperplacefigtwo

\subsection{Throughput Bound and Roofline Interpretation}
Each path imposes a necessary service time on the window: at least $Q_S/\rho$ for streaming and $Q_R/\tau$ for resident writes. The window cannot complete before the longer of the two, and dividing $Q_S$ by this time bounds its throughput:
\begin{align}
  T&\geq\max\!\left(\frac{Q_S}{\rho},\frac{Q_R}{\tau}\right),\label{eq:time}\\
  P&\leq\min(\rho,\tau\SI),\qquad \SI^*=\frac{\rho}{\tau}.
  \label{eq:roofline}
\end{align}
The two service requirements balance at the ridge $\SI^*$. For $\SI<\SI^*$, the resident term is tighter and the window is resident-bound; for $\SI>\SI^*$, the streaming term is tighter and the window is streaming-bound. The bound holds whether or not the two services overlap; shared resources, execution dependencies, and scheduling may keep the actual throughput below it.

Fig.~\ref{fig:roofline-comparison} sets this bound beside the classical Roofline: the demand--throughput pairing is preserved, but work is measured in streaming input bytes rather than operations. Within one execution, $P$ still converts to useful operations per second, $P_{\mathrm{OP}}$: for a logical $128\times128$ INT8 matrix-vector multiplication at two operations per MAC, each input byte corresponds to 256 OP, so $P_{\mathrm{OP}}=(256\ \mathrm{OP}/\Byte)\,P$ for this execution.

\subsection{From Streaming Intensity to Reuse}
The ridge can be expressed as the number of input vectors that one resident load must support. Consider a resident matrix $W[N,K]$ with output dimension $N$ and input dimension $K$, and let the reuse count $U$ be the number of input vectors served after one complete load, including the first use; these vectors may accumulate across batches or requests. With $b_S$ and $b_R$ bytes per streaming and resident element,
\begin{equation}
  Q_S=UKb_S,\qquad Q_R=NKb_R,\qquad \SI=\frac{Ub_S}{Nb_R}.
  \label{eq:matrix-demand}
\end{equation}
At equal byte widths, $\SI=U/N$, as listed for five logical matrices in Table~\ref{tab:reuse-intensity}; for fixed $U$ and precision, $K$ scales both demands equally and cancels.
\paperplacetableone

The same matrix yields different $\SI$ under different residency policies. Reloading before every input fixes $U=1$ per load, however many inputs the window serves. A matrix already resident before the window and never rewritten gives $Q_R=0$, $\SI=\infty$, and $P\leq\rho$; this pre-resident window differs from unbounded reuse after one load, $U\to\infty$, which keeps $Q_R=NKb_R$ finite.

On the hardware side, let $\Delta_S$ be the steady-state service interval of one complete input vector and $T_R$ the service time of one complete load, so that $\rho=Kb_S/\Delta_S$ and $\tau=NKb_R/T_R$. Setting $\SI=\SI^*$ gives the \emph{critical reuse}
\begin{equation}
  U^*=\frac{T_R}{\Delta_S}=\frac{Nb_R}{b_S}\,\SI^*,
  \label{eq:critical-reuse}
\end{equation}
which reduces to $U^*=N\SI^*$ at equal byte widths. The two service requirements thus balance at $U\Delta_S=T_R$: serving $U^*$ input vectors takes as much service time as one complete load. A load-once window is therefore resident-bound for $U<U^*$ and streaming-bound for $U>U^*$. Here, $U$ is the reuse a workload offers, whereas $U^*$ is the reuse the hardware requires, set by its load time relative to its input service interval.

For this load-once reference, normalizing~\eqref{eq:roofline} by $\rho$ gives
\begin{equation}
  \frac{P}{\rho}\leq\min\!\left(1,\frac{\SI}{\SI^*}\right)=\min\!\left(1,\frac{U}{U^*}\right),
  \label{eq:reuse-bound}
\end{equation}
so insufficient reuse caps the throughput at a fraction $U/U^*$ of the streaming ceiling.
